\section{AMC 8 2020}
        \begin{enumerate}
        	\item (\textit{AMC 8 2020}) Luka is making lemonade to sell at a school fundraiser. His recipe requires $4$ times as much water as sugar and twice as much sugar as lemon juice. He uses $3$ cups of lemon juice. How many cups of water does he need? 


 $\textbf{(A) } 6\qquad\textbf{(B) } 8\qquad\textbf{(C) } 12\qquad\textbf{(D) } 18\qquad\textbf{(E) } 24\qquad$


 
 \textbf{Answer: }E



\textbf{Solution 1 (Direct)}

We have $\text{water} : \text{sugar} : \text{lemon juice} = 4\cdot 2 : 2 : 1 = 8 : 2 : 1,$ so Luka needs $3 \cdot 8 = \boxed{\textbf{(E) }24}$ cups.

 



\textbf{Solution 2 (Stepwise)}

Since the amount of sugar is twice the amount of lemon juice, Luka uses $3\cdot2=6$ cups of sugar.

 Since the amount of water is $4$ times the amount of sugar, he uses $6\cdot4=\boxed{\textbf{(E) }24}$ cups of water.

 



\textbf{Solution 3 (Combination of Solutions 1 and 2)}

The ratio is $\text{Water}:\text{Sugar}:\text{Lemon Juice},$ or $8:2:1.$ Since we know that Luka used 3 cups of lemon juice, he needs $3\cdot2=6$ cups of sugar. Because the amount of water is $4$ times the amount of sugar Luka needs, he will need $6\cdot4=\boxed{\textbf{(E) }24}$ cups of water.

 Thanks to MRENTHUSIASM for the inspiration!

 EarthSaver 15:12, 11 June 2021 (EDT)

 


	\item (\textit{AMC 8 2020}) Four friends do yardwork for their neighbors over the weekend, earning $\$15, \$20, \$25,$ and $\$40,$ respectively. They decide to split their earnings equally among themselves. In total, how much will the friend who earned $\$40$ give to the others?


 $\textbf{(A) }\$5 \qquad \textbf{(B) }\$10 \qquad \textbf{(C) }\$15 \qquad \textbf{(D) }\$20 \qquad \textbf{(E) }\$25$


 
 \textbf{Answer: }C



\textbf{Solution}

The friends earn $\$\left(15+20+25+40\right)=\$100$ in total. Since they decided to split their earnings equally, it follows that each person will get $\$\left(\frac{100}{4}\right)=\$25$. Since the friend who earned $\$40$ will need to leave with $\$25$, he will have to give $\$\left(40-25\right)=\boxed{\textbf{(C) }\$15}$ to the others.

 



\textbf{Solution 2}

The friends earn $\$\left(15+20+25+40\right)=\$100$ in total. Since they decided to split their earnings equally, it follows that each person will get $\$\left(\frac{100}{4}\right)=\$25$. The friend who earned $25$ will not distribute anything, the friend that earned $15$ needs $10$, and the friend that earned $20$ needs $5$. $10 + 5 = 15$, therefore the answer is $\$\left(40-25\right)=\boxed{\textbf{(C) }\$15}$.

 


	\item (\textit{AMC 8 2020}) Carrie has a rectangular garden that measures $6$ feet by $8$ feet. She plants the entire garden with strawberry plants. Carrie is able to plant $4$ strawberry plants per square foot, and she harvests an average of $10$ strawberries per plant. How many strawberries can she expect to harvest?
$\textbf{(A) }560 \qquad \textbf{(B) }960 \qquad \textbf{(C) }1120 \qquad \textbf{(D) }1920 \qquad \textbf{(E) }3840$


 
 \textbf{Answer: }D



\textbf{Solution 1}

Note that the unit of the answer is \textbf{strawberries}, which is the product of

 
                \begin{itemize}
                    \item square feet
                \end{itemize}


                
                \begin{itemize}
                    \item plants per square foot
                \end{itemize}


                
                \begin{itemize}
                    \item strawberries per plant
                \end{itemize}


                By conversion factors, we have 
 \[\left(6 \ \color{red}\cancel{\mathrm{ft}}\color{black}\cdot8 \ \color{red}\cancel{\mathrm{ft}}\color{black}\right)\cdot\left(4 \ \frac{\color{green}\cancel{\mathrm{plants}}}{\color{red}\cancel{\mathrm{ft}^2}}\right)\cdot\left(10 \ \frac{\mathrm{strawberries}}{\color{green}\cancel{\mathrm{plant}}}\right)=6\cdot8\cdot4\cdot10 \ \mathrm{strawberries}=\boxed{\textbf{(D) }1920} \ \mathrm{strawberries}.\]

 



\textbf{Solution 2}

The area of the garden is $6 \cdot 8 = 48$ square feet. Since Carrie plants $4$ strawberry plants per square foot, there are a total of $48 \cdot 4=192$ strawberry plants, each of which produces $10$ strawberries on average. Accordingly, she can expect to harvest $192 \cdot 10 = \boxed{\textbf{(D) }1920}$ strawberries.

 


 


	\item (\textit{AMC 8 2020}) Three hexagons of increasing size are shown below. Suppose the dot pattern continues so that each successive hexagon contains one more band of dots. How many dots are in the next hexagonal?


 
                    \begin{figure}[H]
                        \centering
                        \includegraphics[width=0.5\textwidth]{Problem/AMC_8_AMC_2020_4_Problem_Figure_1.png}
                    \end{figure}
                    


 $\textbf{(A) }35 \qquad \textbf{(B) }37 \qquad \textbf{(C) }39 \qquad \textbf{(D) }43 \qquad \textbf{(E) }49$


 
 \textbf{Answer: }B



\textbf{Solution 1 (Pattern of the Rows)}

Looking at the rows of each hexagon, we see that the first hexagon has $1$ dot, the second has $2+3+2$ dots, and the third has $3+4+5+4+3$ dots. Given the way the hexagons are constructed, it is clear that this pattern continues. Hence, the fourth hexagon has $4+5+6+7+6+5+4=\boxed{\textbf{(B) }37}$ dots.

 



\textbf{Solution 2 (Pattern of the Bands)}

The dots in the next hexagon have four bands. From innermost to outermost:

 <ol style="margin-left: 1.5em;"> \item The first band has $1$ dot. 

 \item The second band has $6$ dots: $1$ dot at each vertex of the hexagon. 

 \item The third band has $6+6\cdot1=12$ dots: $1$ dot at each vertex of the hexagon and $1$ other dot on each edge of the hexagon. 

 \item The fourth band has $6+6\cdot2=18$ dots: $1$ dot at each vertex of the hexagon and $2$ other dots on each edge of the hexagon.

 </ol>Together, the answer is $1+6+12+18=\boxed{\textbf{(B) }37}.$

 



\textbf{Solution 3 (Pattern of the Bands)}

The first hexagon has $1$ dot, the second hexagon has $1+6$ dots, the third hexagon has $1+6+12$ dots, and so on. The pattern continues since to go from hexagon $n$ to hexagon $(n+1),$ we add a new band of dots around the outside of the existing ones, with each side of the band having side length $(n+1).$ Thus, the number of dots added is $6(n+1)-6 = 6n$ (we subtract $6$ as each of the corner hexagons in the band is counted as part of two sides.). We therefore predict that the fourth hexagon has $1+6+12+18=\boxed{\textbf{(B) }37}$ dots.

 <u>\textbf{Remark}</u>

 For positive integers $n,$ let $h_n$ denote the number of dots in the $n$th hexagon. We have $h_1=1$ and $h_{n+1}=h_n+6n.$

 It follows that $h_2=7,h_3=19,$ and $h_4=37.$

 



\textbf{Solution 4 (Brute Force)}

From the full diagram below, the answer is $\boxed{\textbf{(B) }37}.$
                    \begin{figure}[H]
                        \centering
                        \includegraphics[width=0.5\textwidth]{Solution/AMC_8_AMC_2020_4_Solution_Figure_1.png}
                    \end{figure}
                    ~MRENTHUSIASM

 



\textbf{Solution 5 (Summation Notation)}

We solve thiss using the $\textbf{Sigma Function}$ (summation notation).

 The sequence starts at:\[a_1 = 1\]

 Each term increases by a growing difference:\[\begin{aligned}a_2 &= a_1 + 6 = 7 \\a_3 &= a_2 + 12 = 19 \\a_4 &= a_3 + 18 = 37\end{aligned}\]

 We observe that the difference added at each step forms an arithmetic sequence:  \[6,\, 12,\, 18,\, 24,\, \dots = 6k \text{ for } k = 1, 2, 3, \dots\]

 To express this with the \textbf{Sigma Function}, we write the $n$-th term as:

 a\_n = 1 + \sum\_{k=1}\textasciicircum{}{n-1} 6k\]

 This is a standard arithmetic sum, so we use the formula:\[\sum_{k=1}^{n-1} k = \frac{(n-1)n}{2}\]

 Substitute this in:\[a_n = 1 + 6 \cdot \frac{(n-1)n}{2}\]

 When we simplify:\[a_n = 1 + 3n(n - 1)\]\[\boxed{a_n = 1 + 3n(n - 1)}\]

 Let’s compute the 4th term using the sigma-based formula:\[a_4 = 1 + 3 \cdot 4 \cdot (4 - 1) = 1 + 3 \cdot 12 = 1 + 36 = \boxed{37}\]So the answer is 37.

 



\textbf{Solution 6 (Quick)}

Notice that the middle row, which is also the row with the most dots in it, has the same number of dots as the topmost or bottommost row multiplied by 2, then minus one. Since the new hexagon will have four dots at the topmost and bottommost rows, the middle row will have (4*2)-1=7 dots. Also, notice that the rows increase in the number of dots from the topmost to the middle row, then decrease from the middle row to the bottommost row by 1. So, noticing these patterns, the amount of dots in the next hexagon is simply $4+5+6+7+6+5+4=\boxed{\textbf{(B) }37}$ dots.~Evžen D.

 



\textbf{Solution 7 (Fast Patern Recognition)}

The first hexagon has $1$ dot in it. The second hexagon has $7$ dots in it. The third hexagon has $19$ dots on it. From the first to the second, we add $6$ hexagons. From the second to the third, we add $12$ hexagons. Therefore, recognizing the pattern, from the third to the fourth case, we have to add $18$ dots to the $19$ we already have, resulting in our answer of $\boxed{\textbf{(B) }37}$ dots.

 


	\item (\textit{AMC 8 2020}) Three fourths of a pitcher is filled with pineapple juice. The pitcher is emptied by pouring an equal amount of juice into each of $5$ cups. What percent of the total capacity of the pitcher did each cup receive? 


 $\textbf{(A) }5 \qquad \textbf{(B) }10 \qquad \textbf{(C) }15 \qquad \textbf{(D) }20 \qquad \textbf{(E) }25$


 
 \textbf{Answer: }C



\textbf{Solution 1}

Each cup is filled with $\frac{3}{4} \cdot \frac{1}{5} = \frac{3}{20}$ of the amount of juice in the pitcher, so the percentage is $\frac{3}{20} \cdot 100 = \boxed{\textbf{(C) }15}$.

 



\textbf{Solution 2}

The pitcher is $\frac{3}{4}$ full, i.e. $75\%$ full. Therefore each cup receives $\frac{75}{5}=\boxed{\textbf{(C) }15}$ percent of the total capacity.

 



\textbf{Solution 3}

Assume that the pitcher has a total capacity of $100$ ounces. Since it is filled three fourths with pineapple juice, it contains $75$ ounces of pineapple juice, which means that each cup will contain $\frac{75}{5}=15$ ounces of pineapple juice. Since the total capacity of the pitcher was $100$ ounces, it follows that each cup received $15\%$ of the total capacity of the pitcher, yielding $\boxed{\textbf{(C) }15}$ as the answer.

 


	\item (\textit{AMC 8 2020}) Aaron, Darren, Karen, Maren, and Sharon rode on a small train that has five cars that seat one person each. Maren sat in the last car. Aaron sat directly behind Sharon. Darren sat in one of the cars in front of Aaron. At least one person sat between Karen and Darren. Who sat in the middle car?


 $\textbf{(A) }\text{Aaron} \qquad \textbf{(B) }\text{Darren} \qquad \textbf{(C) }\text{Karen} \qquad \textbf{(D) }\text{Maren}\qquad \textbf{(E) }\text{Sharon}$


 
 \textbf{Answer: }A



\textbf{Solution 1}

Write the order of the cars as $\square\square\square\square\square$, where the left end of the row represents the back of the train and the right end represents the front. Call the people $A$, $D$, $K$, $M$, and $S$ respectively. The first condition gives $M\square\square\square\square$, so we try $MAS\square\square$, $M\square AS\square$, and $M\square\square AS$. In the first case, as $D$ sat in front of $A$, we must have $MASDK$ or $MASKD$, which do not comply with the last condition. In the second case, we obtain $MKASD$, which works, while the third case is obviously impossible since it results in there being no way for $D$ to sit in front of $A$. It follows that, with the only possible arrangement being $MKASD$, the person sitting in the middle car is $\boxed{\textbf{(A) }\text{Aaron}}$.

 



\textbf{Solution 2}

Follow the first few steps of Solution 1. We must have $M\square\square\square\square$, and also have $AS\text{(anything)}D \text{ and } K\dots D$. There are only $4$ spaces available for $A, S, D, K$, so the only possible arrangement of them is $KASD$, so the arrangement is $MKASD$, so the person in the middle car is $\boxed{\textbf{(A) }\text{Aaron}}$.

 


	\item (\textit{AMC 8 2020}) How many integers between $2020$ and $2400$ have four distinct digits arranged in increasing order? (For example, $2347$ is one integer.)


 $\textbf{(A) }\text{9} \qquad \textbf{(B) }\text{10} \qquad \textbf{(C) }\text{15} \qquad \textbf{(D) }\text{21}\qquad \textbf{(E) }\text{28}$


 
 \textbf{Answer: }C



\textbf{Solution 1}

Firstly, we can observe that the second digit of such a number cannot be $1$ or $2$ because the digits must be distinct and in increasing order. The second digit also cannot be $4$ as the number must be less than $2400$, so the second digit must be $3$. It remains to choose the latter two digits, which must be $2$ distinct digits from $\left\{4,5,6,7,8,9\right\}$. That can be done in $\binom{6}{2} = \frac{6 \cdot 5}{2 \cdot 1} = 15$ ways; there is then only $1$ way to order the digits, namely in increasing order. This means the answer is $\boxed{\textbf{(C) }15}$.

 



\textbf{Solution 2 (without using the "choose" function)}

As in Solution 1, we find that the first two digits must be $23$, and the third digit must be at least $4$ because the digits can not repeat. If it is $4$, then there are $5$ choices for the last digit, namely $5$, $6$, $7$, $8$, or $9$. Similarly, if the third digit is $5$, there are $4$ choices for the last digit, namely $6$, $7$, $8$, and $9$; if $6$, there are $3$ choices; if $7$, there are $2$ choices; and if $8$, there is $1$ choice. It follows that the total number of such integers is $5+4+3+2+1=\boxed{\textbf{(C) }15}$.

 


	\item (\textit{AMC 8 2020}) Ricardo has $2020$ coins, some of which are pennies ($1$-cent coins) and the rest of which are nickels ($5$-cent coins). He has at least one penny and at least one nickel. What is the difference in cents between the greatest possible and least possible amounts of money that Ricardo can have?


 $\textbf{(A) }\text{8062} \qquad \textbf{(B) }\text{8068} \qquad \textbf{(C) }\text{8072} \qquad \textbf{(D) }\text{8076}\qquad \textbf{(E) }\text{8082}$


 
 \textbf{Answer: }C



\textbf{Solution 1}

Clearly, the amount of money Ricardo has will be maximized when he has the maximum number of nickels. Since he must have at least one penny, the greatest number of nickels he can have is $2019$, giving a total of $(2019\cdot 5 + 1)$ cents. Analogously, the amount of money he has will be least when he has the greatest number of pennies; as he must have at least one nickel, the greatest number of pennies he can have is also $2019$, giving him a total of $(2019\cdot 1 + 5)$ cents. Hence the required difference is 
 \[(2019\cdot 5 + 1)-(2019\cdot 1 + 5)=2019\cdot 4-4=4\cdot 2018=\boxed{\textbf{(C) }8072}\]

 



\textbf{Solution 2}

Suppose Ricardo has $p$ pennies, so then he has $(2020-p)$ nickels. In order to have at least one penny and at least one nickel, we require $p \geq 1$ and $2020 - p \geq 1$, i.e. $1 \leq p \leq 2019$. The number of cents he has is $p+5(2020-p) = 10100-4p$, so the maximum is $10100-4 \cdot 1$ and the minimum is $10100 - 4 \cdot 2019$, and the difference is therefore 
 \[(10100 - 4\cdot 1) - (10100 - 4\cdot 2019) = 4\cdot 2019 - 4 = 4\cdot 2018 = \boxed{\textbf{(C) }8072}\]

 


	\item (\textit{AMC 8 2020}) Akash's birthday cake is in the form of a $4 \times 4 \times 4$ inch cube. The cake has icing on the top and the four side faces, and no icing on the bottom. Suppose the cake is cut into $64$ smaller cubes, each measuring $1 \times 1 \times 1$ inch, as shown below. How many small pieces will have icing on exactly two sides?


 
                    \begin{figure}[H]
                        \centering
                        \includegraphics[width=0.5\textwidth]{Problem/AMC_8_AMC_2020_9_Problem_Figure_1.png}
                    \end{figure}
                    


 

$\textbf{(A) }\text{12} \qquad \textbf{(B) }\text{16} \qquad \textbf{(C) }\text{18} \qquad \textbf{(D) }\text{20}\qquad \textbf{(E) }\text{24}$


 
 \textbf{Answer: }D



\textbf{Solution 1}

Notice that, for a small cube which does not form part of the bottom face, it will have exactly $2$ faces with icing on them only if it is one of the $2$ center cubes of an edge of the larger cube. There are $12-4 = 8$ such edges (as we exclude the $4$ edges of the bottom face), so this case yields $2 \cdot 8 = 16$ small cubes. As for the bottom face, we can see that only the $4$ corner cubes have exactly $2$ faces with icing, so the total is $16+4 = \boxed{\textbf{(D) }20}$.

 



\textbf{Solution 2}

The following diagram shows $12$ of the small cubes having exactly $2$ faces with icing on them; that is all of them except for those on the hidden face directly opposite the front face.

 
                    \begin{figure}[H]
                        \centering
                        \includegraphics[width=0.5\textwidth]{Solution/AMC_8_AMC_2020_9_Solution_Figure_1.png}
                    \end{figure}
                    But the hidden face is an exact copy of the front face, so the answer is $12+8=\boxed{\textbf{(D) }20}$.

 



\textbf{Solution 3}

It is observable that only the middle-edged pieces of each face will have the icing exactly on two sides. There are 4 such pieces on each face. Considering 5 faces of the cube (since the bottom is not iced), we can state that the number of such pieces of dimensions 1x1x1 will be 5x4=20. Therefore, our answer is 20.

 



\textbf{Solution 4 (For Rubik's Cubers)}

On a $4$x$4$ Rubik's cube, there are exactly $24$'edge' pieces, $8$ 'corners', and $24$ 'center' pieces. Edge pieces have $2$ frosted faces (the ones on the bottom only have one, corners have $3$ frosted faces, and centers have $1$. So since we have $24$ edges pieces, we subtract the $8$ 'edge' pieces on the bottom (they only have one frosted face), and then we add the $4$ bottom 'corner' pieces (they also have frosted faces). we get $24-8+4=\boxed{\textbf{(D) }20}$.

 -Solution by MismatchedCubing and Andrew\_Lu

 


	\item (\textit{AMC 8 2020}) Zara has a collection of $4$ marbles: an Aggie, a Bumblebee, a Steelie, and a Tiger. She wants to display them in a row on a shelf, but does not want to put the Steelie and the Tiger next to one another. In how many ways can she do this?


 $\textbf{(A) }6 \qquad \textbf{(B) }8 \qquad \textbf{(C) }12 \qquad \textbf{(D) }18 \qquad \textbf{(E) }24$


 
 \textbf{Answer: }C



\textbf{Solution 1}

Let the Aggie, Bumblebee, Steelie, and Tiger, be referred to by $A,B,S,$ and $T$, respectively. If we ignore the constraint that $S$ and $T$ cannot be next to each other, we get a total of $4!=24$ ways to arrange the 4 marbles. We now simply have to subtract out the number of ways that $S$ and $T$ can be next to each other. If we place $S$ and $T$ next to each other in that order, then there are three places that we can place them, namely in the first two slots, in the second two slots, or in the last two slots (i.e. $ST\square\square, \square ST\square, \square\square ST$). However, we could also have placed $S$ and $T$ in the opposite order (i.e. $TS\square\square, \square TS\square, \square\square TS$). Thus there are 6 ways of placing $S$ and $T$ directly next to each other. Next, notice that for each of these placements, we have two open slots for placing $A$ and $B$. Specifically, we can place $A$ in the first open slot and $B$ in the second open slot or switch their order and place $B$ in the first open slot and $A$ in the second open slot. This gives us a total of $6\times 2=12$ ways to place $S$ and $T$ next to each other. Subtracting this from the 

 


	\item (\textit{AMC 8 2020}) After school, Maya and Naomi headed to the beach, $6$ miles away. Maya decided to bike while Naomi took a bus. The graph below shows their journeys, indicating the time and distance traveled. What was the difference, in miles per hour, between Naomi's and Maya's average speeds?


 
                    \begin{figure}[H]
                        \centering
                        \includegraphics[width=0.5\textwidth]{Problem/AMC_8_AMC_2020_11_Problem_Figure_1.png}
                    \end{figure}
                    


 $\textbf{(A) }6 \qquad \textbf{(B) }12 \qquad \textbf{(C) }18 \qquad \textbf{(D) }20 \qquad \textbf{(E) }24$


 
 \textbf{Answer: }E



\textbf{Solution 1}

Naomi travels $6$ miles in a time of $10$ minutes, which is equivalent to $\dfrac{1}{6}$ of an hour. Since $\text{speed} = \frac{\text{distance}}{\text{time}}$, her speed is $\frac{6}{\left(\frac{1}{6}\right)} = 36$ mph. By a similar calculation, Maya's speed is $12$ mph, so the answer is $36-12 = \boxed{\textbf{(E) }24}$.

 



\textbf{Solution 2 (Variant of Solution 1)}

Naomi's speed of $6$ miles in $10$ minutes is equivalent to $6 \cdot 6 = 36$ miles per hour, while Maya's speed of $6$ miles in $30$ minutes (i.e. half an hour) is equivalent to $6 \cdot 2 = 12$ miles per hour. The difference is consequently $36-12=\boxed{\textbf{(E) }24}$.

 


	\item (\textit{AMC 8 2020}) For a positive integer $n,$ the factorial notation $n!$ represents the product of the integers
from $n$ to $1$. (For example, $6! = 6 \cdot 5 \cdot 4 \cdot 3 \cdot 2 \cdot 1$.) What value of $N$ satisfies the following equation?

 \[5! \cdot 9! = 12 \cdot N!\]
$\textbf{(A) }10 \qquad \textbf{(B) }11 \qquad \textbf{(C) }12 \qquad \textbf{(D) }13 \qquad \textbf{(E) }14$


 
 \textbf{Answer: }A



\textbf{Solution 1}

We have $5! = 2 \cdot 3 \cdot 4 \cdot 5$, and $2 \cdot 5 \cdot 9! = 10 \cdot 9! = 10!$. Therefore, the equation becomes $3 \cdot 4 \cdot 10! = 12 \cdot N!$, and so $12 \cdot 10! = 12 \cdot N!$. Cancelling the $12$s, it is clear that $N=\boxed{\textbf{(A) }10}$.

 



\textbf{Solution 2 (variant of Solution 1)}

Since $5! = 120$, we obtain $120\cdot 9!=12\cdot N!$, which becomes $12\cdot 10\cdot 9!=12\cdot N!$ and thus $12 \cdot 10!=12\cdot N!$. We therefore deduce $N=\boxed{\textbf{(A) }10}$.

 



\textbf{Solution 3 (using answer choices and elimination)}

We can see that the answers $\textbf{(B)}$ to $\textbf{(E)}$ contain a factor of $11$, but there is no such factor of $11$ in $5! \cdot 9!$. The factor 11 is in every answer choice except $\boxed{\textbf{(A) }10}$, so four of the possible answers are eliminated. Therefore, the answer must be $\boxed{\textbf{(A) }10}$.~edited by HW73

 



\textbf{Solution 4}

We notice that $5! \cdot 9! = (5!)^2 \cdot (9 \cdot 8 \cdot 7 \cdot 6).$

 We know that $5! = 120,$ so we have $120(5! \cdot 9 \cdot 8 \cdot 7 \cdot 6) = 12 \cdot N!$

 Isolating $N!$ we have $N! = 10 \cdot 5! \cdot 9 \cdot 8 \cdot 7 \cdot 6 \Rightarrow N! = 10! \Rightarrow N = \boxed{\textbf{(A) }10}.$

 



\textbf{Solution 5}

Notice that $9!$ is equivalent to $\frac{12!}{12 \cdot 11 \cdot 10}$. In the expression $\frac{12!}{12 \cdot 11 \cdot 10} \cdot 5 \cdot 4 \cdot 3 \cdot 2$, cancel out $12$, $4$, $3$, and $10$, $5$, $2$. The resulting equation is $\frac{12!}{11} = 12 \cdot N!$. The equation gives $12 \cdot 10! = 12 \cdot N!$. Therefore, N = 10, so the answer is $\boxed{(A)10}$.

 


	\item (\textit{AMC 8 2020}) Jamal has a drawer containing $6$ green socks, $18$ purple socks, and $12$ orange socks. After adding more purple socks, Jamal noticed that there is now a $60\%$ chance that a sock randomly selected from the drawer is purple. How many purple socks did Jamal add?


 $\textbf{(A) }6 \qquad \textbf{(B) }9 \qquad \textbf{(C) }12 \qquad \textbf{(D) }18 \qquad \textbf{(E) }24$


 
 \textbf{Answer: }B



\textbf{Solution 1}

After Jamal adds $x$ purple socks, he has $(18+x)$ purple socks and $6+18+12+x=(36+x)$ total socks. This means the probability of drawing a purple sock is $\frac{18+x}{36+x}$, so we obtain 
 \[\frac{18+x}{36+x}=\frac{3}{5}\] Since $\frac{18+9}{36+9}=\frac{27}{45}=\frac{3}{5}$, the answer is $\boxed{\textbf{(B) }9}$.

 



\textbf{Solution 2 (variant of Solution 1)}

As in Solution 1, we have the equation $\frac{18+x}{36+x}=\frac{3}{5}$. Cross-multiplying yields $90+5x=108+3x \Rightarrow 2x=18 \Rightarrow x=9$. Thus, Jamal added $\boxed{\textbf{(B) }9}$ purple socks.

 



\textbf{Solution 3 very fast}

The answer is $9$

 



\textbf{Solution 4 (Testing Answer Choices)}

The probability that a randomly selected sock is purple is: The \# of purple socks ÷ The total \# of socks. We can test the answer choices until we find the correct answer.

 
[A] If Jamal adds 6 purple socks, he will have 18 + 6 = 24 purple socks. The probability of a randomly selected sock being purple after adding 6 socks is $\frac{24}{6+24+12}=\frac{24}{42}=\frac{4}{7}$, which isn't equal to $\frac{3}{5}$. 

 
[B] If Jamal adds 9 purple socks, he will have 18 + 9 = 27 purple socks. The probability of a randomly selected sock being purple after adding 9 socks is $\frac{27}{6+27+12}=\frac{27}{45}=\frac{3}{5}$, which is indeed equal to $\frac{3}{5}$.

 
So our final answer is $\boxed{\textbf{(B) }9}$ purple socks.

 


	\item (\textit{AMC 8 2020}) There are $20$ cities in the County of Newton. Their populations are shown in the bar chart below. The average population of all the cities is indicated by the horizontal dashed line. Which of the following is closest to the total population of all $20$ cities?


 
                    \begin{figure}[H]
                        \centering
                        \includegraphics[width=0.5\textwidth]{Problem/AMC_8_AMC_2020_14_Problem_Figure_1.png}
                    \end{figure}
                    


 $\textbf{(A) }65{,}000 \qquad \textbf{(B) }75{,}000 \qquad \textbf{(C) }85{,}000 \qquad \textbf{(D) }95{,}000 \qquad \textbf{(E) }105{,}000$


 
 \textbf{Answer: }D



\textbf{Solution}

We can see that the dotted line is exactly halfway between $4{,}500$ and $5{,}000$, so it is at $4{,}750$. As this is the average population of all $20$ cities, the total population is simply $4{,}750 \cdot 20 = \boxed{\textbf{(D) }95{,}000}$.

 



\textbf{Solution 2}

The dotted line is slightly below $5000$. Since $5000\cdot20=100000$, the answer is slightly below this, so $\boxed{\textbf{(D) }95{,}000}$ would be the best choice to guess.

 


	\item (\textit{AMC 8 2020}) Suppose $15\%$ of $x$ equals $20\%$ of $y.$ What percentage of $x$ is $y?$


 $\textbf{(A) }5 \qquad \textbf{(B) }35 \qquad \textbf{(C) }75 \qquad \textbf{(D) }133 \frac13 \qquad \textbf{(E) }300$


 
 \textbf{Answer: }C



\textbf{Solution 1}

Since $20\% = \frac{1}{5}$, multiplying the given condition by $5$ shows that $y$ is $15 \cdot 5 = \boxed{\textbf{(C) }75}$ percent of $x$.

 



\textbf{Solution 2}

Letting $x=100$ (without loss of generality), the condition becomes $0.15\cdot 100 = 0.2\cdot y \Rightarrow 15 = \frac{y}{5} \Rightarrow y=75$. Clearly, it follows that $y$ is $75\%$ of $x$, so the answer is $\boxed{\textbf{(C) }75}$.

 



\textbf{Solution 3}

We have $15\%=\frac{3}{20}$ and $20\%=\frac{1}{5}$, so $\frac{3}{20}x=\frac{1}{5}y$. Solving for $y$, we multiply by $5$ to give $y = \frac{15}{20}x = \frac{3}{4}x$, so the answer is $\boxed{\textbf{(C) }75}$.

 



\textbf{Solution 4}

We are given $0.15x = 0.20y$, so we may assume without loss of generality that $x=20$ and $y=15$. This means $\frac{y}{x}=\frac{15}{20}=\frac{75}{100}$, and thus the answer is $\boxed{\textbf{(C) }75}$.

 



\textbf{Solution 5}

$15\%$ of $x$ is $0.15x$, and $20\%$ of $y$ is $0.20y$. We put $0.15x$ and $0.20y$ into an equation, creating $0.15x = 0.20y$ because $0.15x$ equals $0.20y$. Solving for $y$, dividing $0.2$ to both sides, we get $y = \frac{15}{20}x = \frac{3}{4}x$, so the answer is $\boxed{\textbf{(C) }75}$.

 



\textbf{Solution 6}

$15\%$ of $x$ can be written as $\frac{15}{100}x$, or $\frac{15x}{100}$. $20\%$ of $y$ can similarly be written as $\frac{20}{100}y$, or $\frac{20y}{100}$. So now, $\frac{15x}{100} = \frac{20y}{100}$. Using cross-multiplication, we can simplify the equation as: $1500x = 2000y$. Dividing both sides by $500$, we get: $3x = 4y$. $\frac{3}{4}$ is the same thing as $75\%$, so the answer is $\boxed{\textbf{(C) }75}$.

 


	\item (\textit{AMC 8 2020}) Each of the points $A,B,C,D,E,$ and $F$ in the figure below represents a different digit from $1$ to $6.$ Each of the five lines shown passes through some of these points. The digits along each line are added to produce five sums, one for each line. The total of the five sums is $47.$ What is the digit represented by B?


 
                    \begin{figure}[H]
                        \centering
                        \includegraphics[width=0.5\textwidth]{Problem/AMC_8_AMC_2020_16_Problem_Figure_1.png}
                    \end{figure}
                    


 $\textbf{(A) }1 \qquad \textbf{(B) }2 \qquad \textbf{(C) }3 \qquad \textbf{(D) }4 \qquad \textbf{(E) }5$


 
 \textbf{Answer: }E



\textbf{Solution 1}

We can form the following expressions for the sum along each line:
 \[\begin{dcases}A+B+C\\A+E+F\\C+D+E\\B+D\\B+F\end{dcases}\]Adding these together, we must have $2A+3B+2C+2D+2E+2F=47$, i.e. $2(A+B+C+D+E+F)+B=47$. Since $A,B,C,D,E,F$ are unique integers between $1$ and $6$, we obtain $A+B+C+D+E+F=1+2+3+4+5+6=21$ (where the order doesn't matter as addition is commutative), so our equation simplifies to $42 + B = 47$. This means $B = \boxed{\textbf{(E) }5}$.~RJ5303707

 



\textbf{Solution 2}

Following the first few steps of Solution 1, we have $2(A+C+D+E+F)+3B=47$.      Because an even number $2(A+C+D+E+F)$ subtracted from an odd number (47) is always odd, we know that $3B$ is odd, showing that $B$ is odd. Now we know that $B$ is 1, 3, or 5. If we try $B=1$, we get $43=47$, which is not true. Also when we test $B=3$, we get $45=47$, which is also false. Therefore, we end up getting $B = \boxed{\textbf{(E) }5}$.

 -minor edit by tonykuncheng

 -minor edit by Lebron James

 


	\item (\textit{AMC 8 2020}) How many factors of $2020$ have more than $3$ factors? (As an example, $12$ has $6$ factors, namely $1, 2, 3, 4, 6,$ and $12.$)


 $\textbf{(A) }6 \qquad \textbf{(B) }7 \qquad \textbf{(C) }8 \qquad \textbf{(D) }9 \qquad \textbf{(E) }10$


 
 \textbf{Answer: }B



\textbf{Solution 1}

Since $2020 = 2^2 \cdot 5 \cdot 101$, we can simply list its factors: 
 \[1, 2, 4, 5, 10, 20, 101, 202, 404, 505, 1010, 2020.\] There are $12$ factors; only $1, 2, 4, 5, 101$ don't have over $3$ factors, so the remaining $12-5 = \boxed{\textbf{(B) }7}$ factors have more than $3$ factors.

 



\textbf{Solution 2}

As in Solution 1, we prime factorize $2020$ as $2^2\cdot 5\cdot 101$, and we recall the standard formula that the number of positive factors of an integer is found by adding $1$ to each exponent in its prime factorization, and then multiplying these. Thus $2020$ has $(2+1)(1+1)(1+1) = 12$ factors. The only number which has one factor is $1$. For a number to have exactly two factors, it must be prime, and the only prime factors of $2020$ are $2$, $5$, and $101$. For a number to have three factors, it must be a square of a prime (this follows from the standard formula mentioned above), and from the prime factorization, the only square of a prime that is a factor of $2020$ is $4$. Thus, there are $5$ factors of $2020$ which themselves have $1$, $2$, or $3$ factors (namely $1$, $2$, $4$, $5$, and $101$), so the number of factors of $2020$ that have more than $3$ factors is $12-5=\boxed{\textbf{(B) }7}$.

 



\textbf{Solution 3}

Let $d(n)$ be the number of factors of n. We know by prime factorization that $d(2020) = 12$. These $12$ numbers can be divided into unordered pairs ${a,b}$ where $ab = 2020$. Since $d(2020) = d(a)d(b)$, one of $d(a), d(b)$ has $3$ or less factors and the other has $4$ or more - in to total $6$ factors of $2020$ with more than $3$ factors. However, this argument has exceptions where $a$ and $b$ share a nontrivial common factor, which in this case can only be two. There are two cases - One in which $5$ and $101$ divide the same factor, WLOG assumed to be $a$, so that $d(a) = 2^3 > 3$ and $d(b) = 2$, as otherwise. In the other case, $a = 5\cdot2$ and $b = 101\cdot2$, so that $d(a) = d(b) = 4$. This adds one factor with more than $3$ factors, so the answer is $\boxed{\textbf{(B) }7}$.

 


	\item (\textit{AMC 8 2020}) Rectangle $ABCD$ is inscribed in a semicircle with diameter $\overline{FE},$ as shown in the figure. Let $DA=16,$ and let $FD=AE=9.$ What is the area of $ABCD?$


 
                    \begin{figure}[H]
                        \centering
                        \includegraphics[width=0.5\textwidth]{Problem/AMC_8_AMC_2020_18_Problem_Figure_1.png}
                    \end{figure}
                    


 $\textbf{(A) }240 \qquad \textbf{(B) }248 \qquad \textbf{(C) }256 \qquad \textbf{(D) }264 \qquad \textbf{(E) }272$


 
 \textbf{Answer: }A



\textbf{Solution 1 (Pythagorean Theorem)}


                    \begin{figure}[H]
                        \centering
                        \includegraphics[width=0.5\textwidth]{Solution/AMC_8_AMC_2020_18_Solution_Figure_1.png}
                    \end{figure}
                    

 Let $O$ be the center of the semicircle. The diameter of the semicircle is $9+16+9=34$, so $OC = 17$. By symmetry, $O$ is the midpoint of $DA$, so $OD=OA=\frac{16}{2}= 8$. By the Pythagorean theorem in right-angled triangle $ODC$ (or $OBA$), we have that $CD$ (or $AB$) is $\sqrt{17^2-8^2}=15$. Accordingly, the area of $ABCD$ is $16\cdot 15=\boxed{\textbf{(A) }240}$.

 



\textbf{Solution 2 (Coordinate Geometry)}

Let the midpoint of segment $FE$ be the origin. Evidently, point $D=(-8,0)$ and $A=(8,0)$. Since points $C$ and $B$ share $x$-coordinates with $D$ and $A$ respectively, it suffices to find the $y$-coordinate of $B$ (which will be the height of the rectangle) and multiply this by $DA$ (which we know is $16$). The radius of the semicircle is $\frac{9+16+9}{2} = 17$, so the whole circle has equation $x^2+y^2=289$; as already stated, $B$ has the same $x$-coordinate as $A$, i.e. $8$, so substituting this into the equation shows that $y=\pm15$. Since $y>0$ at $B$, the y-coordinate of $B$ is $15$. Therefore, the answer is $16\cdot 15 = \boxed{\textbf{(A) }240}$.

 (Note that the synthetic solution (Solution 1) is definitely faster and more elegant. However, this is the solution that you should use if you can't see any other easier strategy.)

 



\textbf{Solution 3}

We can use a result from the Art of Problem Solving \textit{Introduction to Algebra} book Sidenote: for a semicircle with diameter $(1+n)$, such that the $1$ part is on one side and the $n$ part is on the other side, the height from the end of the $1$ side (or the start of the $n$ side) is $\sqrt{n}$. To use this formula, we scale the figure down by $9$; this will give the height a length of $\sqrt{\frac{16+9}{9}} = \sqrt{\frac{25}{9}} = \frac{5}{3}$. Now, scaling back up by $9$, the height $DC$ is $9 \cdot \frac{5}{3} = 15$. The answer is then $15 \cdot 16 = \boxed{\textbf{(A) }240}$.-\href{https://artofproblemsolving.com/wiki/index.php?title=User:Sweetmango77}{SweetMango77}

 



\textbf{Solution 4 (Power of a Point Theorem)}

Draw the other half of the circle as follows:
                    \begin{figure}[H]
                        \centering
                        \includegraphics[width=0.5\textwidth]{Solution/AMC_8_AMC_2020_18_Solution_Figure_2.png}
                    \end{figure}
                    By the \href{https://artofproblemsolving.com/wiki/index.php?title=Power\_of\_a\_Point\_Theorem}{Power of a Point Theorem}, $FD\cdot DE = CD\cdot C'D$. By symmetry, $CD = C'D$. We see that $FD = 9$ and $DE = 16 + 9 = 25$. Substituting in these values, $9\cdot 25 = CD^2$, giving $CD^2 = 225$ and $CD = 15$.The area of the rectangle is therefore $15\cdot 16 = \boxed{\textbf{(A) }240}$.

 



\textbf{Solution 5 (Vertical Theorem)}


                    \begin{figure}[H]
                        \centering
                        \includegraphics[width=0.5\textwidth]{Solution/AMC_8_AMC_2020_18_Solution_Figure_3.png}
                    \end{figure}
                    

 According to the Pythagorean Theorem and the Vertical Theorem, we can find out that $OG=\sqrt{\left(\frac{2\times9+16}{2}\right)^2 - \left(\frac{16}{2}\right)^2}=15$. Therefore, the answer is $15\times16=\boxed{\textbf{(A) }240}$;

 


	\item (\textit{AMC 8 2020}) A number is called flippy if its digits alternate between two distinct digits. For example, $2020$ and $37373$ are flippy, but $3883$ and $123123$ are not. How many five-digit flippy numbers are divisible by $15?$


 $\textbf{(A) }3 \qquad \textbf{(B) }4 \qquad \textbf{(C) }5 \qquad \textbf{(D) }6 \qquad \textbf{(E) }8$


 
 \textbf{Answer: }B



\textbf{Solution 1}

A number is divisible by $15$ precisely if it is divisible by $3$ and $5$. The latter means the last digit must be either $5$ or $0$, and the former means the sum of the digits must be divisible by $3$. If the last digit is $0$, the first digit would be $0$ (because the digits alternate), which is not possible. Hence the last digit must be $5$, and the number is of the form $5\square 5\square 5$. If the unknown digit is $x$, we deduce $5+x+5+x+5 \equiv 0 \pmod{3} \Rightarrow 2x \equiv 0 \pmod{3}$. We know $2^{-1}$ exists modulo $3$ because 2 is relatively prime to 3, so we conclude that $x$ (i.e. the second and fourth digit of the number) must be a multiple of $3$. It can be $0$, $3$, $6$, or $9$, so there are $\boxed{\textbf{(B) }4}$ options: $50505$, $53535$, $56565$, and $59595$.

 



\textbf{Solution 2 (variant of Solution 1, no mods)}

As in Solution 1, we find that such numbers must start with $5$ and alternate with $5$ (i.e. must be of the form $5\square 5\square 5$), where the two digits between the $5$s need to be the same. Call that digit $x$. For the number to be divisible by $3$, the sum of the digits must be divisible by $3$; since the sum of the three $5$s is $15$, which is already a multiple of $3$, it must also be the case that $x+x=2x$ is a multiple of $3$. Thus, the problem reduces to finding the number of digits from $0$ to $9$ for which $2x$ is a multiple of $3$. This leads to $x=0$, $3$, $6$, or $9$, so there are $\boxed{\textbf{(B) }4}$ possible numbers (namely $50505$, $53535$, $56565$, and $59595$).

 



\textbf{Solution 3}

After finding out that the last digit must be $5$, the number is of the form $5\square 5\square 5$. If the unknown digit is $x$, we can find that one of the solutions to $x$ is $0$, since $5+5+5$ is equal to $15$, which is divisible by $3$. After trying every one digit number, you'll notice that $x$ must be a multiple of $3$, meaning that $x=0$, $3$, $6$, or $9$. $50505$, $53535$, $56565$, and $59595$ are the $\boxed{\textbf{(B) }4}$ solutions to this question.

 



\textbf{Solution 4 (mods)}

assume the number is $ababa$$10101a+1010b=0 (mod 15)\newline$$6a+5b=0 (mod 15)\newline$$a=0 (mod 5)\newline$$5b=0 (mod 15)\newline$$b=0 (mod 3)\newline$Solutions: $(5,0),(5,3),(5,6),(5,9)\newline$$\boxed{4}$

 


	\item (\textit{AMC 8 2020}) A scientist walking through a forest recorded as integers the heights of $5$ trees standing in a row. She observed that each tree was either 2 times as tall or half as tall as the one to its right. Unfortunately some of her data was lost when rain fell on her notebook. Her notes are shown below, with blanks indicating the missing numbers. Based on her observations, the scientist was able to reconstruct the lost data. What was the average height of the trees, in meters?


 
 \[\begingroup \setlength{\tabcolsep}{10pt} \renewcommand{\arraystretch}{1.5} \begin{tabular}{|c|c|} \hline Tree 1 & \rule{0.4cm}{0.15mm} meters \\ Tree 2 & 11 meters \\ Tree 3 & \rule{0.5cm}{0.15mm} meters \\ Tree 4 & \rule{0.5cm}{0.15mm} meters \\ Tree 5 & \rule{0.5cm}{0.15mm} meters \\ \hline Average height & \rule{0.5cm}{0.15mm}\text{ .}2 meters \\ \hline \end{tabular} \endgroup\]
$\textbf{(A) }22.2 \qquad \textbf{(B) }24.2 \qquad \textbf{(C) }33.2 \qquad \textbf{(D) }35.2 \qquad \textbf{(E) }37.2$


 
 \textbf{Answer: }B



\textbf{Solution 1}

We will show that $22$, $11$, $22$, $44$, and $22$ meters are the heights of the trees from left to right. We are given that all tree heights are integers, so since Tree 2 has height $11$ meters, we can deduce that Trees 1 and 3 both have a height of $22$ meters. There are now three possible cases for the heights of Trees 4 and 5 (in order for them to be integers), namely heights of $11$ and $22$, $44$ and $88$, or $44$ and $22$. Checking each of these, in the first case, the average is $17.6$ meters, which doesn't end in $.2$ as the problem requires. Therefore, we consider the other cases. With $44$ and $88$, the average is $37.4$ meters, which again does not end in $.2$, but with $44$ and $22$, the average is $24.2$ meters, which does. Consequently, the answer is $\boxed{\textbf{(B) }24.2}$.

 



\textbf{Solution 2: EASY TO UNDERSTAND!!!}

Notice the average height of the trees ends with $0.2$; therefore, the sum of all five heights of the trees must end with $1$ or $6$. ($0.2$ $\cdot$ $5$ = $1$) We already know Tree $2$ is $11$ meters tall. Both Tree $1$ and Tree $3$ must $22$ meters tall - since neither can be $5.5$ because according to the problem it says “integers”. Once again, apply our observation for solving for the Tree $4$'s height. Tree $4$ can't be $11$ meters for the sum of the five tree heights to still end with $1$. Therefore, the Tree $4$ is $44$ meters tall. Now, Tree $5$ can either be $22$ or $88$. Find the average height for both cases of Tree $5$. Doing this, we realize the Tree $5$ must be $22$ for the average height to end with $0.2$ and that the average height is $\boxed{\textbf{(B)}\ 24.2}$.

 -by DummyAMC

 



\textbf{Solution 3}

As in Solution 1, we shall show that the heights of the trees are $22$, $11$, $22$, $44$, and $22$ meters. Let $S$ be the sum of the heights, so that the average height will be $\frac{S}{5}$ meters. We note that $0.2 = \frac{1}{5}$, so in order for $\frac{S}{5}$ to end in $.2$, $S$ must be one more than a multiple of $5$. Moreover, as all the heights are integers, the heights of Tree 1 and Tree 3 are both $22$ meters. At this point, our table looks as follows:
 \[\begingroup \setlength{\tabcolsep}{10pt} \renewcommand{\arraystretch}{1.5} \begin{tabular}{|c|c|} \hline Tree 1 & 22 meters \\ Tree 2 & 11 meters \\ Tree 3 & 22 meters \\ Tree 4 & \rule{0.5cm}{0.15mm} meters \\ Tree 5 & \rule{0.5cm}{0.15mm} meters \\ \hline Average height & \rule{0.5cm}{0.15mm}\text{ .}2 meters \\ \hline \end{tabular} \endgroup\]

 If Tree 4 now has a height of $11$, then Tree 5 would need to have height $22$, but in that case $S$ would equal $88$, which is not $1$ more than a multiple of $5$. So we instead take Tree 4 to have height $44$. Then the sum of the heights of the first 4 trees is $22+11+22+44 = 99$, so using a height of $22$ for Tree 5 gives $S=121$, which is $1$ more than a multiple of $5$ (whereas $88$ gives $S = 187$, which is not). Thus the average height of the trees is $\frac{121}{5} = \boxed{\textbf{(B) }24.2}$ meters.

 



\textbf{Solution 4}

Since we know that the tree heights have to be integers, then it is immediate that Tree 1 and 3 have a height of $22$. Now using the information given by the last column (that the average of the heights of the trees ends in $.2$), we can tell that the sum of all the heights of the trees ends in either $1$ or $6$, because those are the only numbers from $0$ to $9$ that are congruent to $1$ after taking modulo $5$. The two multiples of eleven (eleven because all of the tree heights have to be a multiple of eleven if they are integers) that come to mind are $66$ and $121$. Since the sum of the heights of Tree 1, 2, and 3 is already $55$, we know that $66$ is impossible to obtain. Then, we can decide with relative confidence that the answer should be $\frac{121}{5} = \boxed{\textbf{(B) }24.2}$.

 
P.S. : We can check our solution by trying to obtain $121$. Surely enough, we have Tree 1 = $22$, Tree 2 = $11$, Tree 3 = $22$, Tree 4 = $44$, and Tree 5 = $22$. These all add up to get $121$.

 



\textbf{Solution 5 (basic modular arithmetic)}

Since her average ends with 0.2, the sum of her tree heights must be $\equiv 1 \pmod {5}$.If Tree 1's height is $a$, Tree 3's is $b$, Tree 4's is $c$, and Tree 5's is $d$, then we $a$ and $b$ are both 22, since all of the tree heights are integers.Now we have $22+11+22+c+d \equiv 1 \pmod {5}$. Simplifying, we get $c+d \equiv 1 \pmod {5}$.The only possible combination of $c$ and $d$ that abides the condition are 44 and 22, respectively.Adding these up, we get 121, and $\frac{121}{5}$ is $\boxed{\textbf{(B) }24.2}$.

 -kn07

 


	\item (\textit{AMC 8 2020}) A game board consists of $64$ squares that alternate in color between black and white. The figure below shows square $P$ in the bottom row and square $Q$ in the top row. A marker is placed at $P.$ A step consists of moving the marker onto one of the adjoining white squares in the row above. How many $7$-step paths are there from $P$ to $Q?$ (The figure shows a sample path.)


 
                    \begin{figure}[H]
                        \centering
                        \includegraphics[width=0.5\textwidth]{Problem/AMC_8_AMC_2020_21_Problem_Figure_1.png}
                    \end{figure}
                    


 $\textbf{(A) }28 \qquad \textbf{(B) }30 \qquad \textbf{(C) }32 \qquad \textbf{(D) }33 \qquad \textbf{(E) }35$


 
 \textbf{Answer: }A



\textbf{Solution 1}

Notice that, from one of the $1$ or $2$ white squares immediately beneath it (since the marker can only move on white squares). This means that the number of ways to move from $P$ to that square is the sum of the number of ways to move from $P$ to each of the white squares immediately beneath it (also called the Waterfall Method). To solve the problem, we can accordingly construct the following diagram, where each number in a square is calculated as the sum of the numbers on the white squares immediately beneath that square (and thus will represent the number of ways to remove from $P$ to that square, as already stated).

 
                    \begin{figure}[H]
                        \centering
                        \includegraphics[width=0.5\textwidth]{Solution/AMC_8_AMC_2020_21_Solution_Figure_1.png}
                    \end{figure}
                    

 The answer is therefore $\boxed{\textbf{(A) }28}$.

 Note: This is a classic example of a problem involving \href{https://artofproblemsolving.com/wiki/index.php?title=Pascal\%27s\_triangle}{Pascal's triangle}.

 



\textbf{Solution 2}

Suppose we "extend" the chessboard infinitely with $2$ additional columns to the right, as shown below. The red line shows the right-hand edge of the original board.

 
                    \begin{figure}[H]
                        \centering
                        \includegraphics[width=0.5\textwidth]{Solution/AMC_8_AMC_2020_21_Solution_Figure_2.png}
                    \end{figure}
                    

 The total number of paths from $P$ to $Q$, including invalid paths which cross over the red line, is then the number of paths which make $4$ steps up-and-right and $3$ steps up-and-left, which is $\binom{4+3}{3} = \binom{7}{3} = 35$. We need to subtract the number of invalid paths, i.e. the number of paths that pass through $X$ or $Y$. To get to $X$, the marker has to make $3$ up-and-right steps, after which it can proceed to $Q$ with $3$ steps up-and-left and $1$ step up-and-right. Thus, the number of paths from $P$ to $Q$ that pass through $X$ is $1 \cdot \binom{3+1}{3} = 4$. Similarly, the number of paths that pass through $Y$ is $\binom{4+1}{1}\cdot 1 = 5$. However, we have now double-counted the invalid paths which pass through both $X$ and $Y$; from the diagram, it is clear that there are only $2$ of these (as the marker can get from $X$ to $Y$ by a step up-and-left and a step-up-and-right in either order). Hence the number of invalid paths is $4+5-2=7$, and the number of valid paths from $P$ to $Q$ is $35-7 = \boxed{\textbf{(A) }28}$.

 



\textbf{Solution 3}

On any white square, we may choose to go left or right, as long as we do not cross over the border of the board. Call the moves $L$ and $R$ respectively. Every single legal path consists of $4$ $R's$ and $3$ $L's$, so now all we have to find is the number of ways to order $4 R's$ and $3 L's$ in any way, which is ${7 \choose 3}=35$. However, we originally promised that we will not go over the border, and now we have to subtract the paths that do go over the border. The paths that go over the border are any paths that start with RRR(1 path), RR(5 paths) and LRRR(1 path) so our final number of paths is $35-7=\boxed{(A)=28}.$~PEKKA

 



\textbf{Solution 4: Intuitive Approach}

We label the rows starting from the bottom. At row 1, there is $1$ way: at P. We draw all the possible ways to get to Q. There are two ways to choose for row 2, and another two ways to choose for row 3. However, you can go to the "edge" or the farthest possible square westward of Q, so you can't multiply by 2 again. Notice how, at the first step, we figured that the answer was even, so choice D and E are eliminated, and after the second row, we realized it must be a multiple of 4, so choice B is eliminated. When we get to the fourth row, we do not multiply by 2 again, since we have limited possibilities rather than multiplying by 2 again. Choice C implies that there are two possibilities per row; however, we know that if you go to the farthest possible, you only have one possibility, so it is not $2^5 = 32$ so we know that the answer is choice $\boxed{\textbf{(A)}}$. ~hh99754539

 


	\item (\textit{AMC 8 2020}) When a positive integer $N$ is fed into a machine, the output is a number calculated according to the rule shown below.


 
                    \begin{figure}[H]
                        \centering
                        \includegraphics[width=0.5\textwidth]{Problem/AMC_8_AMC_2020_22_Problem_Figure_1.png}
                    \end{figure}
                    


 For example, starting with an input of $N=7,$ the machine will output $3 \cdot 7 +1 = 22.$ Then if the output is repeatedly inserted into the machine five more times, the final output is $26.$

 \[7 \to 22 \to 11 \to 34 \to 17 \to 52 \to 26\]When the same $6$-step process is applied to a different starting value of $N,$ the final output is $1.$ What is the sum of all such integers $N?$

 \[N \to \rule{0.5cm}{0.15mm} \to \rule{0.5cm}{0.15mm} \to \rule{0.5cm}{0.15mm} \to \rule{0.5cm}{0.15mm} \to \rule{0.5cm}{0.15mm} \to 1\]


 $\textbf{(A) }73 \qquad \textbf{(B) }74 \qquad \textbf{(C) }75 \qquad \textbf{(D) }82 \qquad \textbf{(E) }83$


 
 \textbf{Answer: }E



\textbf{Solution 1}

We start with final output of $1$ and work backward, taking cares to consider all possible inputs that could have resulted in any particular output. This produces following set of possibilities each stage:
 \[\{1\}\rightarrow\{2\}\rightarrow\{4\}\rightarrow\{1,8\}\rightarrow\{2,16\}\rightarrow\{4,5,32\}\rightarrow\{1,8,10,64\}\]where, for example, $2$ must come from $4$ (as there is no integer $n$ satisfying $3n+1=2$), but $16$ could come from $32$ or $5$ (as $\frac{32}{2} = 3 \cdot 5 + 1 = 16$, and $32$ is even while $5$ is odd). By construction, last set in this sequence contains all the numbers which will lead to number $1$ to end of the $6$-step process, and sum is $1+8+10+64=\boxed{\textbf{(E) }83}$.

 



\textbf{Solution 2 (variant of Solution 1)}

As in Solution 1, we work backwards from $1$, this time showing the possible cases in a tree diagram:

 
                    \begin{figure}[H]
                        \centering
                        \includegraphics[width=0.5\textwidth]{Solution/AMC_8_AMC_2020_22_Solution_Figure_1.png}
                    \end{figure}
                    

 The possible numbers are those at the "leaves" of the tree (the ends of the various branches), which are $1$, $8$, $64$, and $10$. Thus the answer is $1+8+64+10=\boxed{\textbf{(E) }83}$.

 



\textbf{Solution 3 (algebraic)}

We begin by finding the inverse of the function that the machine uses. Call the input $I$ and the output $O$. If $I$ is even, $O=\frac{I}{2}$, and if $I$ is odd, $O=3I+1$. We can therefore see that $I=2O$ when $I$ is even and $I=\frac{O-1}{3}$ when $I$ is odd. Therefore, starting with $1$, if $I$ is even, $I=2$, and if $I$ is odd, $I=0$, but the latter is not valid since $0$ is not actually odd. This means that the 2nd-to-last number in the sequence has to be $2$. Now, substituting $2$ into the inverse formulae, if $I$ is even, $I=4$ (which is indeed even), and if $I$ is odd, $I=\frac{1}{3}$, which is not an integer. This means the 3rd-to-last number in the sequence has to be $4$. Substituting in $4$, if $I$ is even, $I=8$, but if $I$ is odd, $I=1$. Both of these are valid solutions, so the 4th-to-last number can be either $1$ or $8$. If it is $1$, then by the argument we have just made, the 5th-to-last number has to be $2$, the 6th-to-last number has to be $4$, and the 7th-to-last number, which is the first number, must be either $1$ or $8$. In this way, we have ultimately found two solutions: $N=1$ and $N=8$.

 On the other hand, if the 4th-to-last number is $8$, substituting $8$ into the inverse formulae shows that the 5th-to-last number is either $16$ or $\frac{7}{3}$, but the latter is not an integer. Substituting $16$ shows that if $I$ is even, $I=32$, but if I is odd, $I=5$, and both of these are valid. If the 6th-to-last number is $32$, then the first number must be $64$, since $\frac{31}{3}$ is not an integer; if the 6th-to-last number is $5,$ then the first number has to be $10$, as $\frac{4}{3}$ is not an integer. This means that, in total, there are $4$ solutions for $N$, specifically, $1$, $8$, $10$, and $64$, which sum to $\boxed{\textbf{(E) }83}$.

 \textbf{Remark} This function is the basis for the \href{https://en.wikipedia.org/wiki/Collatz\_conjecture}{Collatz conjecture}, which asserts that every positive integer $n$ will eventually output the cycle 1, 4, 2, 1, 4, 2, ... when this function is repeatedly iterated starting with $n$. This conjecture remains unsolved but has been tested for all starting integers up to $2^{71}$

 


	\item (\textit{AMC 8 2020}) Five different awards are to be given to three students. Each student will receive at least one award. In how many different ways can the awards be distributed? 


 $\textbf{(A) }120 \qquad \textbf{(B) }150 \qquad \textbf{(C) }180 \qquad \textbf{(D) }210 \qquad \textbf{(E) }240$


 
 \textbf{Answer: }B



\textbf{Solution 1 (Constructive Counting)}

Firstly, observe that a single student can't receive $4$ or $5$ awards because this would mean that one of the other students receives no awards. Thus, each student must receive either $1$, $2$, or $3$ awards. If a student receives $3$ awards, then the other two students must each receive $1$ award; if a student receives $2$ awards, then another student must also receive $2$ awards and the remaining student must receive $1$ award. We consider each of these two cases in turn.

 If a student receives three awards, there are $3$ ways to choose which student this is, and $\binom{5}{3}$ ways to give that student $3$ out of the $5$ awards. Next, there are $2$ students left and $2$ awards to give out, with each student getting one award. There are clearly just $2$ ways to distribute these two awards out, giving $3\cdot\binom{5}{3}\cdot 2=60$ ways to distribute the awards in this case.

 In the other case, two students receive $2$ awards and one student receives $1$ award. We know there are $3$ choices for which student gets $1$ award. There are $\binom{3}{1}$ ways to do this. Then, there are $\binom{5}{2}$ ways to give the first student his two awards, leaving $3$ awards yet to be distributed. There are then $\binom{3}{2}$ ways to give the second student his $2$ awards. Finally, there is only $1$ student and $1$ award left, so there is only $1$ way to distribute this award. This results in $\binom{5}{2}\cdot\binom{3}{2}\cdot 1\cdot 3 =90$ ways to distribute the awards in this case. Adding the results of these two cases, we get $60+90=\boxed{\textbf{(B) }150}$.

 



\textbf{Solution 2 (Version 2 of Solution 1)}

Upon inspection (specified in the above solution), there are two cases of the distribution of awards to the students: one student gets 3 awards and the other each get 1 award or one student gets 1 award and the other two get 2 awards. 

 
In the first case, there are $\binom{3}{1} = 3$ ways to choose the person who gets 3 awards. From here, there are $\binom{5}{3} = 10$ ways to choose the 3 awards from the 5 total awards. Now, one person has $2$ choices for awards and the other has $1$ choice for the award. Thus, the total number of ways to choose awards in this case is $3 \cdot 10 \cdot 2 \cdot 1 = 60$.  

 
In the other case, there are $\binom{3}{1} = 3$ ways to choose the person who gets 1 award, and $5$ choices for his/her award. Then, one person has $\binom{4}{2} = 6$ ways to have his/her awards and the other person has $\dbinom{2}{2} = 1$ ways to have his/her awards. This gives $3 \cdot 5 \cdot 6 \cdot 1  = 90$ ways for this case. 

 Adding these numbers together, we get $60 + 90 = 150$ ways to distribute the awards, or choice $\boxed{\textbf{(B) }150}$.

 



\textbf{Solution 3 (Complementary Counting)}

Without the restriction that each student receives at least one award, we could take each of the awards and choose one of the $3$ students to give it to. This would be $3^5$ ways to distribute the awards in total. Now we need to subtract the cases where at least one student doesn't receive an award. If a student doesn't receive an award, there are $3$ choices for which student that is, so $2^5$ ways of choosing a student to receive each of the awards; in total, $3\cdot32=96$.

 However, if $2$ students both don't receive an award, then this case would be counted twice among the $96$, so we need to add back in these cases. Said in other words, $2$ students not receiving an award is equivalent to $1$ student receiving $5$ awards, and there are $3$ choices for whom that student would be. To finish, the total number of ways to distribute the awards is $243 - 96+3$, or $\boxed{\textbf{(B) }150}$.

 


	\item (\textit{AMC 8 2020}) A large square region is paved with $n^2$ gray square tiles, each measuring $s$ inches on a side. A border $d$ inches wide surrounds each tile. The figure below shows the case for $n=3$. When $n=24$
, the $576$ gray tiles cover $64\%$ of the area of the large square region. What is the ratio $\frac{d}{s}$ for this larger value of $n?$


 
                    \begin{figure}[H]
                        \centering
                        \includegraphics[width=0.5\textwidth]{Problem/AMC_8_AMC_2020_24_Problem_Figure_1.png}
                    \end{figure}
                    


 $\textbf{(A) }\frac6{25} \qquad \textbf{(B) }\frac14 \qquad \textbf{(C) }\frac9{25} \qquad \textbf{(D) }\frac7{16} \qquad \textbf{(E) }\frac9{16}$


 
 \textbf{Answer: }A



\textbf{Solution 1}

The area of the shaded region is $(24s)^2$. To find the area of the large square, we note that there is a $d$-inch border between each of the $23$ pairs of consecutive squares, as well as from between the first/last squares and the large square, for a total of $23+2 = 25$ times the length of the border, i.e. $25d$. Adding this to the total length of the consecutive squares, which is $24s$, the side length of the large square is $(24s+25d)$, yielding the equation $\frac{(24s)^2}{(24s+25d)^2}=\frac{64}{100}$. Taking the square root of both sides (and using the fact that lengths are non-negative) gives $\frac{24s}{24s+25d}=\frac{8}{10} = \frac{4}{5}$, and cross-multiplying now gives $120s = 96s + 100d \Rightarrow 24s = 100d \Rightarrow \frac{d}{s} = \frac{24}{100} = \boxed{\textbf{(A) }\frac{6}{25}}$.

 
Note: Once we obtain $\frac{24s}{24s+25d} = \frac{4}{5},$ to ease computation, we may take the reciprocal of both sides to yield $\frac{24s+25d}{24s} = 1 + \frac{25d}{24s} = \frac{5}{4},$ so $\frac{25d}{24s} = \frac{1}{4}.$ Multiplying both sides by $\frac{24}{25}$ yields the same answer as before. 

 



\textbf{Solution 2}

WLOG (Without Loss of Generality), we may let $s=1$ (since $d$ will be determined by the scale of $s$, and we are only interested in the ratio $\frac{d}{s}$). Then, as the total area of the $576$ gray tiles is simply $576$, the large square has area $\frac{576}{0.64} = 900$, making the side of the large square $\sqrt{900}=30$. As in Solution 1, the side length of the large square consists of the total length of the gray tiles and $25$ lots of the border, so the length of the border is $d = \frac{30-24}{25} = \frac{6}{25}$. Since $\frac{d}{s}=d$ if $s=1$, the answer is $\boxed{\textbf{(A) }\frac{6}{25}}$.

 



\textbf{Solution 3 (using answer choices)}

As in Solution 2, we let $s = 1$ without loss of generality. For sufficiently large $n$, we can approximate the percentage of the area covered by the gray tiles by subdividing most of the region into congruent squares, as shown:
                    \begin{figure}[H]
                        \centering
                        \includegraphics[width=0.5\textwidth]{Solution/AMC_8_AMC_2020_24_Solution_Figure_1.png}
                    \end{figure}
                    	Each red square has side length $(1+d)$, so by solving $\frac{1^2}{(1+d)^2} = \frac{64}{100} \iff \frac{1}{1+d} = \frac{4}{5}$, we obtain $d = \frac{1}{4}$. The actual fraction of the total area covered by the gray tiles will be slightly less than $\frac{1}{(1+d)^2}$, which implies $\frac{1}{(1+d)^2} > \frac{64}{100} \iff \frac{1}{1+d} > \frac{4}{5} \iff d < \frac{1}{4}$. Hence $d$ (and thus $\frac{d}{s}$, since we are assuming $s=1$) is less than $\frac{1}{4}$, and the only choice that satisfies this is $\boxed{\textbf{(A) }\frac{6}{25}}$.

 Note the answer cannot be $\frac{1}{4}$ because this is only an approximation. The closest answer choice is  A.

 



\textbf{Solution 4}

Pack all tiles together to form a square.  As this square $(24s \times 24s)$ covers $64\%$ of the big square, that means its length is $80\%$ long of the big square's $(24s+25d \times 24s+25d)$. Thus, $25d : 24s = 2 : 8$ implies $\frac{d}{s} = \boxed{\textbf{(A) }\frac{6}{25}}$.

 
                    \begin{figure}[H]
                        \centering
                        \includegraphics[width=0.5\textwidth]{Solution/AMC_8_AMC_2020_24_Solution_Figure_2.png}
                    \end{figure}
                    

 


	\item (\textit{AMC 8 2020}) Rectangles $R_1$ and $R_2,$ and squares $S_1,\,S_2,\,$ and $S_3,$ shown below, combine to form a rectangle that is 3322 units wide and 2020 units high. What is the side length of $S_2$ in units?


 
                    \begin{figure}[H]
                        \centering
                        \includegraphics[width=0.5\textwidth]{Problem/AMC_8_AMC_2020_25_Problem_Figure_1.png}
                    \end{figure}
                    


 $\textbf{(A) }651 \qquad \textbf{(B) }655 \qquad \textbf{(C) }656 \qquad \textbf{(D) }662 \qquad \textbf{(E) }666$


 
 \textbf{Answer: }A



\textbf{Solution 1}

Let the side length of each square $S_k$ be $s_k$. Then, from the diagram, we can line up the top horizontal lengths of $S_1$, $S_2$, and $S_3$ to cover the top side of the large rectangle, so $s_{1}+s_{2}+s_{3}=3322$. Similarly, the short side of $R_2$ will be $s_1-s_2$, and lining this up with the left side of $S_3$ to cover the vertical side of the large rectangle gives $s_{1}-s_{2}+s_{3}=2020$. We subtract the second equation from the first to obtain $2s_{2}=1302$, and thus $s_{2}=\boxed{\textbf{(A) }651}$.

 



\textbf{Solution 2}

Assuming that the problem is well-posed, it should be true in the particular case where $S_1 \cong S_3$ and $R_1 \cong R_2$. Let the sum of the side lengths of $S_1$ and $S_3$ be $x$, and let the length of square $S_2$ be $y$. We then have the system 
 \[\begin{dcases}x+y =3322 \\x-y=2020\end{dcases}\] which we solve to determine $y=\boxed{\textbf{(A) }651}$.

 



\textbf{Solution 3 (faster version of Solution 1)}

Since, for each pair of rectangles, the side lengths have a sum of $3322$ or $2020$ and a difference of $S_2$, the answer must be $\dfrac{3322 - 2020}{2} = \dfrac{1302}{2} = \boxed{\textbf{(A) }651}$.

 



\textbf{Solution 4}

Let the side length of $S_2$ be s, and the shorter side length of $R_1$ and $R_2$ be $r$. We have

 
                    \begin{figure}[H]
                        \centering
                        \includegraphics[width=0.5\textwidth]{Solution/AMC_8_AMC_2020_25_Solution_Figure_1.png}
                    \end{figure}
                    

 From this diagram, it is evident that $r+s+r=2020$. Also, the side length of $S_1$ and $S_3$ is $r+s$. Then, $r+s+s+r+s=3322$. Now, we have 2 systems of equations.

 
 \begin{align*}r+s+r &= 2020 \\ r+s+r+s+s &= 3322 \\ \end{align*}
 

 We can see an $r+s+r$ in the 2nd equation, so substituting that in gives us $2020+2s=3322 \Rightarrow 2s= 1302 \Rightarrow s=\boxed{\textbf{(A) }651}$.

 



\textbf{Solution 5 (Cheese)}

Note that the problem is underspecified. We can modify the diagram in any way as long as the relative positioning of the internal rectangles and squares is maintained. Here's an animation describing all such possible configurations:

 
                    \begin{figure}[H]
                        \centering
                        \includegraphics[width=0.5\textwidth]{Solution/AMC_8_AMC_2020_25_Solution_Figure_2.gif}
                    \end{figure}
                    Let's pick the configuration where S1 is as large as possible and takes up the whole height of the rectangle. Then our diagram is just:

 
                    \begin{figure}[H]
                        \centering
                        \includegraphics[width=0.5\textwidth]{Solution/AMC_8_AMC_2020_25_Solution_Figure_3.png}
                    \end{figure}
                    

 For \( S_2 \) and \( S_3 \) to share a top edge, they must actually be congruent! Therefore, the width of the entire rectangle, being given as 3322, is also \( 2020 + 2x \). Solving for \( x \), we get $x = \boxed{\textbf{(A) } 651}$.

 


\end{enumerate}